\section{Algorithmic Proof of the Path-Coloring Lemma}

Let $G = (V, E)$ be the directed graph we intend to color, satisfying the degree
and 2-cycle restrictions outlined previously.

\subsection{Phase 1: Constructing the Initial Coloring}

The first phase aims to partition the edges of $G$ into two color classes such
that no node acts as a junction for two edges of the same color pointing in the
same direction.

We construct an auxiliary undirected graph $A_1 = (E, F_1)$. Every edge $e \in
E$ of the original graph $G$ becomes a vertex in $A_1$ (referred to as an
\textit{A-node}). The edge set $F_1$ (containing \textit{A-edges}) is
constructed to capture structural conflicts: an undirected edge $\{e, f\} \in
F_1$ exists if and only if the directed edges $e$ and $f$ share either a common
head or a common tail in $G$. Specifically, $\{e, f\} \in F_1$ if there exist
vertices $u, v, w \in V$ such that either $e = (u, v)$ and $f = (w, v)$, or $e =
(v, u)$ and $f = (v, w)$. This guarantees that colliding edges are forced into
separate color classes.

Because the maximum total degree of any vertex in $G$ is bounded by 3, every
A-node in $A_1$ has a degree of at most 2. As a result, the auxiliary graph
$A_1$ decomposes entirely into a collection of isolated vertices, simple paths
(A-paths), and simple cycles (A-cycles).

Crucially, every A-cycle in $A_1$ must contain an even number of A-edges. This
property arises because traversing an A-cycle corresponds to alternating between
"head-meets" and "tail-meets" in the original graph $G$. This structural
guarantee allows us to trivially 2-color the A-nodes by simply alternating
colors along each A-path and A-cycle.

\begin{lemma} \label{lemma:initial_coloring}
The edges of $G$ can be colored with two colors such that each color class forms
a collection of node-disjoint paths and cycles. Furthermore, this initial
coloring can be computed in polynomial time.
\end{lemma}
While Lemma \ref{lemma:initial_coloring} successfully isolates paths and cycles,
the path-coloring property requires the complete elimination of all
monochromatic cycles. Note that for any A-path or A-cycle in $A_1$, there are
exactly two valid alternating colorings, and we initially choose one
arbitrarily. This flexibility is the key observation.

\subsection{Phase 2: Refining the Coloring and Breaking Cycles}
To systematically eliminate monochromatic cycles, we classify the A-nodes (edges
of $G$) based on their configuration in $A_1$:
\begin{itemize}
    \item Let $Z$ be the set of all A-nodes belonging to A-cycles in $A_1$.
    \item Let $P$ be the set of all A-nodes belonging to A-paths that contain at
    least one A-edge (i.e., non-isolated A-nodes not in $Z$).
    \item Let $S$ be the set of all isolated A-nodes in $A_1$.
\end{itemize}
Furthermore, for any A-path, we define its first and last A-nodes as
\textit{A-end nodes}.
We first restrict our focus to the subgraph $G' = (V, Z \cup P)$. Our objective
is to ensure that $G'$ can be 2-colored such that each color class contains
strictly node-disjoint paths. Because $G$ lacks 2-cycles, any monochromatic
cycle formed in $G'$ must consist of at least three edges. More importantly, due
to the established alternating coloring, the only way a monochromatic cycle can
manifest in $G'$ is if all participating edges are A-end nodes of various
A-paths in $P$. Furthermore, degree restrictions and the fact that every A-path
in $P$ represents at least two edges in $G$ guarantee that any such
monochromatic cycles in $G'$ are entirely node-disjoint.

To destroy a monochromatic cycle containing edges $e_1, e_2, \dots, e_l$ in
$G'$, we forcefully introduce a color conflict between two adjacent edges. We
arbitrarily select $e_1$ and $e_2$ and add a new auxiliary edge $\{e_1, e_2\}$
to our conflict set. Applying this to all monochromatic cycles in $G'$ produces
an augmented edge set $F_2$, giving a new auxiliary graph $A_2 = (Z \cup P,
F_2)$.

Since the problematic cycles in $G'$ are node-disjoint, there are at most
$|E|/3$ such cycles, meaning they can be identified and neutralized in
polynomial time. While introducing the new A-edges in $F_2$ may fuse some
A-paths into new A-cycles, these newly formed A-cycles maintain an even length.
This is because adding an A-edge $\{e_1, e_2\}$ structurally represents the
meeting of two heads or two tails, preserving the alternating parity required
for 2-colorability.

\begin{lemma} \label{lemma:refined_coloring}
The edges of the subgraph $G' = (V, Z \cup P)$ can be 2-colored such that each
color class consists strictly of node-disjoint paths. This refined coloring is
computable in polynomial time.
\end{lemma}

Following this augmentation, we update the sets $Z$ and $P$ to reflect the
transfer of A-nodes from paths into the newly formed A-cycles of $A_2$. The
final step of the algorithm must then correctly integrate the isolated edges of
set $S$ without closing any new monochromatic cycles.

\subsection{Phase 3: Integrating Isolated Edges}
\paragraph{The Special Case}
Assume temporarily that any directed cycle in $G$ contains at least two edges
belonging to $S$. Under this assumption, we construct a new auxiliary directed
multigraph, $B$, designed to track the monochromatic paths already established
in the refined subgraph $G' = (V, Z \cup P)$:
\begin{itemize}
    \item \textbf{Vertices:} The vertices of $B$ correspond directly to the
    edges in $S$.
    \item \textbf{Edges:} A directed A-edge $(e, f)$ is added to $B$ if there
    exists a monochromatic path in $G'$ connecting the head of $e$ to the tail
    of $f$. This new A-edge inherits the color of that connecting path.
    \item \textbf{Adjacency Handling:} If two edges in $S$ are directly adjacent
    in $G$ (e.g., $e = (u, v)$ and $f = (v, w)$), two parallel A-edges are added
    between them in $B$, one for each color.
\end{itemize}

Because of our temporary assumption that every cycle in $G$ contains at least
two edges from $S$, the multigraph $B$ is guaranteed to be free of self-loops.
Furthermore, degree constraints dictate that every vertex in $B$ has an indegree
and outdegree of at most 2, and any incoming or outgoing pairs of A-edges must
be assigned different colors.

\begin{lemma} \label{lemma:b_coloring}
If the vertices of $B$ can be 2-colored such that every monochromatic cycle in
$B$ contains at least one vertex of the opposite color, then the entirety of $G$
can be validly 2-colored into node-disjoint paths.
\end{lemma}

To deterministically find the coloring required by Lemma \ref{lemma:b_coloring},
we analyze the monochromatic cycles of $B$ by constructing a secondary
undirected multigraph, $B^*$. The vertices of $B^*$ represent the monochromatic
cycles of $B$, and an edge exists between two vertices in $B^*$ if their
corresponding cycles share a vertex in $B$.

The coloring algorithm for $B$ proceeds recursively based on the concept of
\textit{free A-nodes} (vertices in $B$ that belong to at most one monochromatic
cycle):
\begin{enumerate}
    \item If a connected component in $B$ contains a free A-node, we assign it
    the color opposite to its cycle, virtually destroying that cycle. This
    removal breaks the component into smaller subcomponents, each guaranteed to
    now possess a free A-node, allowing the process to recurs.
    \item If a component lacks a free A-node, its corresponding structure in
    $B^*$ must contain a cycle or a double edge. Because all cycles in $B^*$
    have even length, we can alternatively 2-color the elements of the $B^*$
    cycle. This simultaneous color assignment shatters all associated cycles in
    $B$, generating free A-nodes and returning us to the first case.
\end{enumerate}

\paragraph{The General Case: Inductive Graph Reduction}
To remove the restrictive assumption that every cycle in $G$ contains at least
two edges from $S$, the algorithm uses induction on the number of vertices,
$|V|$.

First, trivially isolated edges in $S$ (those sharing no endpoints with any
other edge in $E$) are colored arbitrarily, and adjacent pairs in $S$ are
assigned opposite colors and removed from consideration. Let $T \subseteq S$ be
the remaining subset of these isolated A-nodes.

For each edge $t \in T$, we determine if its tail is reachable from its head
within the currently colored subgraph $(V, Z \cup P)$. If no such path exists
for any $t$, then any cycle formed by incorporating $T$ must necessarily contain
at least two edges from $T$. This directly satisfies our earlier assumption,
allowing us to safely apply the coloring logic of multigraph $B$.

However, if there exists a cycle containing exactly one edge $e \in T$, the
algorithm executes a structural graph contraction. Because $G$ contains no
2-cycles, $e$ must be adjacent to an edge $f$ that is an endpoint of an
established path. By strategically discarding $e$, $f$, and their shared
junction, and rewiring the surrounding edges to bypass the removed segment, we
produce a strictly smaller graph $\hat{G}$. This smaller graph inherits the
necessary structural properties, allowing the path-coloring to proceed
inductively.

\paragraph{Inductive Step and Subcase Analysis}
Throughout this induction, the two colors are denoted $0$ and $1$, so that for a
color $C \in \{0,1\}$ the opposite color is $1-C$.

The discard-and-rewire step splits into several subcases, all worked out by
Bl\"aser \cite{Blaser2004} except one. Below we number them Case 1 through Case
3: Cases 1 and 2 are representative examples of case types Bl\"aser already
covers in full (the remaining case types follow the same pattern and are left to
\cite{Blaser2004}), while Case 3 is the one genuinely missing from his figures.

Figure~\ref{fig:pc-full} shows the local picture we start from: the edge
$e=(x,y)$ we are trying to remove, its preceding path-edge $f_1$, the edge
continuing the cycle at $y$ (call it $g_1$), and the other edge at $x$, $a_1$,
whose far endpoint we write $u$. The other edge at $y$ is $c_1$, with far
endpoint $q$.

\begin{figure}[H]
\centering
\begin{tikzpicture}[>=Stealth, dot/.style={circle,fill=black,inner sep=1.6pt}, hol/.style={circle,draw,fill=white,inner sep=1.6pt}]
    \node[dot,label=above:{$z$}] (z) at (0,0) {};
    \node[dot,label=above:{$x$}] (x) at (1.6,0) {};
    \node[dot,label=above:{$y$}] (y) at (3.2,0) {};
    \node[dot,label=above:{$w_1$}] (w1) at (4.9,0) {};
    \node[hol,label=below:{$u$}] (u) at (1.6,-1.4) {};
    \node[hol,label=below:{$q$}] (q) at (3.2,-1.4) {};
    \draw[->] (z) -- (x) node[midway,above] {$f_1$};
    \draw[->] (x) -- (y) node[midway,above] {$e$};
    \draw[->] (y) -- (w1) node[midway,above] {$g_1$};
    \draw[->] (u) -- (x) node[midway,left] {$a_1$};
    \draw[->] (y) -- (q) node[midway,right] {$c_1$};
\end{tikzpicture}
\caption{Local picture around $e$, with all four surrounding edges
$f_1,e,a_1,c_1$ shown, before anything is discarded.}
\label{fig:pc-full}
\end{figure}

Discarding $e,f_1$ removes $x$ and $z$ from consideration and rewires $a_1$ to
end at $y$ instead of $x$, giving Figure~\ref{fig:pc-branch}. Everything that
follows is based on comparing $u$ (where $a_1$ now starts) with $q$ (where $c_1$
ends).

\begin{figure}[H]
\centering
\begin{tikzpicture}[>=Stealth, dot/.style={circle,fill=black,inner sep=1.6pt}, hol/.style={circle,draw,fill=white,inner sep=1.6pt}]
    \node[dot,label=above:{$y$}] (y) at (1.6,0) {};
    \node[dot,label=above:{$w_1$}] (w1) at (3.3,0) {};
    \node[hol,label=below:{$u$}] (u) at (0.4,-1.4) {};
    \node[hol,label=below:{$q$}] (q) at (2.4,-1.4) {};
    \draw[->] (y) -- (w1) node[midway,above] {$g_1$};
    \draw[->] (u) -- (y) node[midway,left] {$a_1$};
    \draw[->] (y) -- (q) node[midway,right] {$c_1$};
\end{tikzpicture}
\caption{After discarding $e,f_1$ and rewiring $a_1$. Whether $u=q$ decides the
case.}
\label{fig:pc-branch}
\end{figure}

\subparagraph{Case 1 ($u \ne q$): nothing else collides.}
No 2-cycle is created, so we recurse on Figure~\ref{fig:pc-branch} directly,
with $a_1$ ending at $y$ as drawn. Let $C$ be the color $a_1$ receives from the
recursion. We then set $e=C$ and $f_1=1-C$ (Figure~\ref{fig:pc-case0}).

\begin{figure}[H]
\centering
\begin{tikzpicture}[>=Stealth, dot/.style={circle,fill=black,inner sep=1.6pt}, hol/.style={circle,draw,fill=white,inner sep=1.6pt}]
    \node[dot,label=above:{$z$}] (z) at (0,0) {};
    \node[dot,label=above:{$x$}] (x) at (1.6,0) {};
    \node[dot,label=above:{$y$}] (y) at (3.2,0) {};
    \node[dot,label=above:{$w_1$}] (w1) at (4.9,0) {};
    \node[hol,label=below:{$u$}] (u) at (1.6,-1.4) {};
    \node[hol,label=below:{$q$}] (q) at (3.2,-1.4) {};
    \draw[dashed,->] (z) -- (x) node[midway,above] {$f_1$};
    \draw[->] (x) -- (y) node[midway,above] {$e$};
    \draw[->] (u) -- (x) node[midway,left] {$a_1$};
    \draw[dotted,->] (y) -- (w1) node[midway,above] {$g_1$};
    \draw[dotted,->] (y) -- (q) node[midway,right] {$c_1$};
\end{tikzpicture}
\caption{Recovered coloring for $u\ne q$, undoing the reduction in full back to
Figure~\ref{fig:pc-full}: $a_1,e$ share color $C$ (solid) and $f_1$ gets $1-C$
(dashed), matching \texttt{path\_coloring.py}. $g_1,c_1$ (dotted) already have
real colors too, assigned by this same recursive call -- just not ones this rule
determines, so we leave them unstyled.}
\label{fig:pc-case0}
\end{figure}

\subparagraph{Case 2 ($u=q=v$, and $v$ is the tail of the third edge $h$ at
$v$)}
Beyond $a_1,c_1$ colliding, nothing else needs to be touched.

\begin{figure}[H]
\centering
\begin{tikzpicture}[>=Stealth, dot/.style={circle,fill=black,inner sep=1.6pt}, hol/.style={circle,draw,fill=white,inner sep=1.6pt}]
    \node[dot,label=above:{$y$}] (y) at (0,0) {};
    \node[dot,label=left:{$v$}] (v) at (0,-1.4) {};
    \node[dot,label=above:{$w_1$}] (w1) at (2.2,0) {};
    \node[hol] (p) at (2.2,-1.4) {};
    \draw[->] (y) -- (w1) node[midway,above] {$g_1$};
    \draw[->] (v) -- (p) node[midway,below] {$h$};
    \draw[->] (v) to[bend right=35] node[right,xshift=2pt] {$a_1$} (y);
    \draw[->] (y) to[bend right=35] node[left,xshift=-2pt] {$c_1$} (v);
\end{tikzpicture}
\caption{Setup: $a_1,c_1$ close a 2-cycle at $v$, while $h$ leaves $v$ toward an
unrelated node.}
\end{figure}

\begin{figure}[H]
\centering
\begin{tikzpicture}[>=Stealth, dot/.style={circle,fill=black,inner sep=1.6pt}, hol/.style={circle,draw,fill=white,inner sep=1.6pt}]
    \node[dot,label=left:{$y{=}v$}] (yv) at (0,-0.7) {};
    \node[dot,label=above:{$w_1$}] (w1) at (2.2,0) {};
    \node[hol] (p) at (2.2,-1.4) {};
    \draw[->] (yv) -- (w1) node[midway,above] {$g_1$};
    \draw[->] (yv) -- (p) node[midway,below] {$h$};
\end{tikzpicture}
\caption{Discard $a_1,c_1$ and identify $y$ with $v$. Recurse on the smaller
graph.}
\end{figure}
Undoing this reduction in full means restoring not just $a_1,c_1$ but also $x,z$
with $e,f_1$. Crucially, $a_1$ reverts to its original form $(v,x)$ -- the
rewired $(v,y)$ version existed only inside the reduction. We set $a_1$ to
$g_1$'s color, $c_1$ to $h$'s color, then $e$ to $a_1$'s color and $f_1$ to the
opposite.
\begin{figure}[H]
\centering
\begin{tikzpicture}[>=Stealth, dot/.style={circle,fill=black,inner sep=1.6pt}, hol/.style={circle,draw,fill=white,inner sep=1.6pt}]
    \node[dot,label=above:{$z$}] (z) at (0,0) {};
    \node[dot,label=above:{$x$}] (x) at (1.6,0) {};
    \node[dot,label=above:{$y$}] (y) at (3.2,0) {};
    \node[dot,label=above:{$w_1$}] (w1) at (5.4,0) {};
    \node[dot,label=below:{$v$}] (v) at (3.2,-1.4) {};
    \node[hol] (p) at (5.4,-1.4) {};
    \draw[->] (z) -- (x) node[midway,above] {$f_1$};
    \draw[dashed,->] (x) -- (y) node[midway,above] {$e$};
    \draw[dashed,->] (y) -- (w1) node[midway,above] {$g_1$};
    \draw[->] (v) -- (p) node[midway,below] {$h$};
    \draw[dashed,->] (v) -- (x) node[midway,below left] {$a_1$};
    \draw[->] (y) -- (v) node[midway,right] {$c_1$};
\end{tikzpicture}
\caption{Recovered coloring, undone in full: $g_1,a_1,e$ share one color
(dashed), $h,c_1,f_1$ the other (solid). The dashed class is the path $v\to x\to
y\to w_1$ and the solid class is $z\to x$ together with $y\to v\to p$ -- both
simple, no cycle.}
\end{figure}

\subparagraph{Case 3 ($u=q=v$, and $v$ is the head of $h$): omitted from
\cite{Blaser2004}.}
Write $w_1=\textit{head}(g_1)$ and $w_2=\textit{tail}(h)$. If $w_1=w_2$ this is
a loop, not a fresh 2-cycle, fixed by discarding $g_1,c_1,h,e,f_1$ and
recoloring from the next chain edge at $w_1$. Assume $w_1 \ne w_2$. Now $w_1$
carries a second edge $c_2$ (its chain neighbor $g_2$ leads elsewhere, dashed
below) -- and it matters whether $c_2$ happens to land exactly on $w_2$.

\begin{figure}[H]
\centering
\begin{tikzpicture}[>=Stealth, dot/.style={circle,fill=black,inner sep=1.6pt}, hol/.style={circle,draw,fill=white,inner sep=1.6pt}]
    \node[dot,label=above:{$y$}] (y) at (0,0) {};
    \node[dot,label=left:{$v$}] (v) at (0,-1.4) {};
    \node[dot,label=above:{$w_1$}] (w1) at (2.2,0) {};
    \node[dot,label=below:{$w_2$}] (w2) at (2.2,-1.4) {};
    \node[hol] (chain) at (4,0) {};
    \draw[->] (y) -- (w1) node[midway,above] {$g_1$};
    \draw[->] (v) to[bend right=35] node[right,xshift=2pt] {$a_1$} (y);
    \draw[->] (y) to[bend right=35] node[left,xshift=-2pt] {$c_1$} (v);
    \draw[->] (w2) -- (v) node[midway,below] {$h$};
    \draw[->,dashed] (w1) -- (chain) node[midway,above] {\small $g_2$};
    \draw[->,line width=1.1pt] (w1) -- (w2) node[midway,right] {$c_2$};
\end{tikzpicture}
\caption{Setup: $a_1,c_1$ close a 2-cycle at $v$, and $w_1$'s loose edge $c_2$
(bold) happens to coincide with $w_2=\textit{tail}(h)$ -- unrelated to where the
chain edge $g_2$ (dashed) actually goes.}
\end{figure}

\begin{figure}[H]
\centering
\begin{tikzpicture}[>=Stealth, dot/.style={circle,fill=black,inner sep=1.6pt}]
    \node[dot,label=above:{$w_1$}] (w1) at (0,0) {};
    \node[dot,label=above:{$w_2$}] (w2) at (2.2,0) {};
    \draw[->,line width=1.1pt] (w1) -- (w2) node[midway,above] {$c_2$};
\end{tikzpicture}
\caption{Discard $g_1,h,c_1,a_1$ and the nodes $x,y,v$, leaving $c_2$ untouched.
Recurse on what remains.}
\end{figure}

Let $v_e$ be the third edge at $w_2$ (other than $h,c_2$), colored $C$ by the
recursion. We set
\begin{equation}
f_1 = C, \quad e = C, \quad c_1 = C, \qquad g_1 = 1-C, \quad h = 1-C, \quad a_1 = 1-C,
\end{equation}
and leave $c_2$ with whatever color the recursion already gave it.

\begin{figure}[H]
\centering
\begin{tikzpicture}[>=Stealth, dot/.style={circle,fill=black,inner sep=1.6pt}]
    \node[dot,label=above:{$z$}] (z) at (0,0) {};
    \node[dot,label=above:{$x$}] (x) at (1.6,0) {};
    \node[dot,label=above:{$y$}] (y) at (3.2,0) {};
    \node[dot,label=below:{$v$}] (v) at (3.2,-1.4) {};
    \node[dot,label=above:{$w_1$}] (w1) at (5.4,0) {};
    \node[dot,label=below:{$w_2$}] (w2) at (5.4,-1.4) {};
    \draw[->] (z) -- (x) node[midway,above] {$f_1$};
    \draw[->] (x) -- (y) node[midway,above] {$e$};
    \draw[dashed,->] (y) -- (w1) node[midway,above] {$g_1$};
    \draw[dashed,->] (v) -- (x) node[midway,below left] {$a_1$};
    \draw[->] (y) -- (v) node[midway,right] {$c_1$};
    \draw[dashed,->] (w2) -- (v) node[midway,below] {$h$};
    \draw[->,line width=1.1pt] (w1) -- (w2) node[midway,right] {$c_2$};
\end{tikzpicture}
\caption{Recovered coloring, undone in full back to $z,x$, with $a_1$ reverted
to $(v,x)$: $f_1,e,c_1$ (color $C$, solid) versus $g_1,a_1,h$ (color $1-C$,
dashed). The edge $c_2$ (bold) keeps whatever color the recursion assigned it.}
\end{figure}

Bl\"aser's own Fig.~A.10 in \cite{Blaser2004} handles a different coincidence:
the chain edge $g_2$ (not $c_2$) leading to a node whose \emph{own} further edge
(his $c_3$) coincides with $h$, resolved by discarding $g_2,g_1,h,c_1,a_1$ and
splicing a further chain edge $g_3$ into $w_1$. The case above -- $c_2$ itself
landing on $w_2$ -- is not among his figures. It is, however, exactly the branch
\texttt{two\_cycle\_edge == c\_2} in the reference implementation
(\texttt{path\_coloring.py}, function
\texttt{\_lemma\_15\_inductive\_coloring}), which is why we have spelled it out
in full.
\section{Complexity and Running Time}
The efficiency of this coloring scheme relies fundamentally on the graph's
degree restrictions. Since every vertex has a total degree of at most 3, the
number of edges is strictly bounded by $|E| \le \frac{3}{2}|V|$. The inductive
reduction eliminates at least one vertex per step, so the reduction is applied
at most $|V|$ times. Each individual step, however, does not run in $O(|E|)$
time: constructing the auxiliary graph $A_1$ requires comparing every pair of
edges for a shared endpoint, which takes $O(|E|^2) = O(|V|^2)$ time, and
locating the target edge $e$ requires testing, for each of the $O(|V|)$
candidates in $T$, whether it lies on a cycle of $G'$, at a further cost of
$O(|V|+|E|) = O(|V|)$ per candidate -- again $O(|V|^2)$ in total. Hence each
inductive step costs $O(|V|^2)$ time, and since at most $|V|$ steps are
performed, the entire path-coloring algorithm executes deterministically in
$O(|V|) \cdot O(|V|^2) = O(n^3)$ time.
